

%قبل از شروع مسئله ترجیحاً چیزی ننویسید. اما اگر خواستید بنویسید، در قالب لیستی (فقرات و شمارش) بنویسید.
%برای اضافه کردن نمره، بعدی از شروع‌مسئله، [x] را اضافه کنید و به جای x  نمره مورد نظر را قرار دهید. برای زیرمسئله‌ها هم به همین صورت است.
%\شروع‌مسئله[10]
\شروع‌مسئله []
%\متن‌سیاه{سوال عملی اول:}\\
%\متن‌سیاه{دراپ‌تست}\\

\متن‌لاتین{n} فرد داریم که هر کدام یک قدرت مشخص دارند. در ابتدا هر فرد در گروه تک‌نفره‌ی خودش است. تنوع یک گروه برابر است با ضرب ماکسیمم قدرمطلق اختلاف قدرت بین اعضای گروه در تعداد جفت‌های متمایز از اعضای گروه که ماکسیمم اختلاف قدرت را دارند. به طور مثال اگر قدرت اعضای گروهی ۲ و ۲ و ۳ و ۴ باشد تنوع این گروه برابر ۴ = (۲ - ۴) * ۲ است.\\
حال در هر پرسش از شما خواسته می‌شود یکی از عملیات‌های زیر را انجام دهید:\\
۱. \متن‌لاتین{1 a b} : یعنی گروه فرد a و b را با هم ادغام کنید.\\
۲. \متن‌لاتین{2 a} : تنوع گروه فرد a را چاپ کنید.\\
توجه کنید که در درخواست نوع اول گروه a و b لزوماً متفاوت نیستند. در این صورت تغییری در گروه‌ها رخ نمی‌دهد.

\متن‌سیاه{ورودی:}\\
در خط اول \متن‌لاتین{n} تعداد کل افراد آمده است.\\
خط بعدی شامل  \متن‌لاتین{n} عدد صحیح است که \متن‌لاتین{s[i]} قدرت فرد i مشخص می‌کند.\\
در خط بعد  \متن‌لاتین{q} آمده که تعداد عملیات‌هاست.\\ 
در هر یک از q خط بعدی مطابق یکی از عملیات‌ها سه یا دو عدد آمده است.

\متن‌سیاه{خروجی:}\\
به ازای هر کدام از درخواست‌های نوع اول تنوع گروه خواسته شده را چاپ کنید.

\متن‌سیاه{محدودیت‌ها:}\\
$1 \leq n, q \leq 10^{6}$\\
$0 \leq |s[i]| \leq 10^{9}$\\
$1 \leq a, b \leq n$\\

\متن‌سیاه{ورودی و خروجی نمونه:}\\

\begin{latin}
	\begin{tabularx}{\textwidth}{|X | X|}
		\hline
		\textbf{stdin} &
		\textbf{stdout}\\
		\hline
		5 & 0\\
		1 11 4 1 4 & 6\\
		8 & 12\\
		2 1 & 20\\
		1 1 3 & \\
		1 1 5&\\
		2 3&\\
		1 4 5&\\
		2 1&\\
		1 1 2&\\
		2 2&\\
		\hline
	\end{tabularx}
\end{latin}

%}
% صورت اصلی مسئله را در این‌جا بنویسید.
%اگر مسئله چند قسمت داشت به این صورت عمل کنید:
%\شروع{زیرمسائل}
%\زیرمسئله[0] % مسئلۀ ۱آ
%\زیرمسئله % مسئلۀ ۱ب
%\زیرمسئله % مسئلۀ ۱پ
%\زیرمسئله % مسئلۀ ۱ت
%\پایان{زیرمسائل}



%%%%%%%%%%%%%%%%%%%%%%%%%%%%%%%%%%%%%%%%%%%%%%%%%%%%%
% برخی دستورات مفید برای تمرین‌ها. %
%%%%%%%%%%%%%%%%%%%%%%%%%%%%%%%%%%%%%%%%%%%%%%%%%%%%%

% فرمول‌های نمونه:
%$f(x) = \Theta(x)$
%$T(x, y) \leq \log(x) + 2^y + \binom{2n}{n}$
% {\tt code\_function}


% لیست موارد بدون شماره‌گذاری:
%\شروع{فقرات}
%    \فقره مورد اول 
%        \شروع{فقرات}
%            \فقره مورد اول داخلی 
%                \شروع{فقرات}
%                    \فقره اولین مورد خیلی داخلی 
%                    \فقره دومین مورد خیلی داخلی 
%                \پایان{فقرات}
%            \فقره مورد دوم داخلی 
%        \پایان{فقرات}
%    \فقره مورد دوم 
%\پایان{فقرات}

% لیست موارد شماره‌گذاری شده:
%\شروع{شمارش}
%    \فقره یک مورد 
%    \فقره یک مورد دیگر 
%    \فقره باز هم یک مورد دیگر
%\پایان{شمارش}

% برای نوشتن یک فرمول چند خطی که یک شماره بخوره. برای حالت‌های دیگر به سایت زیر مراجعه کنید:
% https://www.overleaf.com/learn/latex/Aligning_equations_with_amsmath
%\begin{align} 
%    \begin{split}
%        (x+y)^3 &= (x+y)^2(x+y) \\
%                &= (x^2+2xy+y^2)(x+y) \\
%                &= (x^3+2x^2y+xy^2) + (x^2y+2xy^2+y^3) \\
%                &= x^3+3x^2y+3xy^2+y^3
%    \end{split}                                 
%\end{align}

% نوشتن ماتریس و آرایه دوبعدی
%\begin{align}
%    A = 
%    \begin{bmatrix}
%        A_{11} & A_{21} \\
%        A_{21} & A_{22}
%    \end{bmatrix}
%\end{align}

% اضافه کردن تصویر با شرح. برای اضافه کردن تصویر بدون شرح، فقط دستور درج‌تصویر را استفاده کنید.
%\شروع{شکل}[h!]
%\تنظیم‌ازوسط
%  \درج‌تصویر[پهنا=0.5\پهنای‌متن]{img.jpg}
%\شرح{شرح تصویر}
%\برچسب{تصویر:۱}
%\پایان{شکل}
%\پایان{زیرمسائل}

%کد زیر را ببینید.
%%\lstinputlisting{source_filename.py}
%\begin{latin}
%\begin{lstlisting}
%# Import the modules
%import sys
%import random
%
%ans = True
%
%while ans:
%question = raw_input("Ask the magic 8 ball a question: (press enter to quit) ")
%
%answers = random.randint(1,8)
%
%if question == "":
%sys.exit()
%
%elif answers == 1:
%print "It is certain"
%
%elif answers == 2:
%print "Outlook good"
%
%elif answers == 3:
%print "You may rely on it"
%
%elif answers == 4:
%print "Ask again later"
%
%elif answers == 5:
%print "Concentrate and ask again"
%
%elif answers == 6:
%print "Reply hazy, try again"
%
%elif answers == 7:
%print "My reply is no"
%
%elif answers == 8:
%print "My sources say no"
%\end{lstlisting}
%\end{latin}

	
	